%! AP Statistics — Unit 4, Lesson 4.2 — Lesson Plan
\documentclass[10pt]{article}
\usepackage{apstats-article}
\usepackage{apstats-boxes}
\usepackage{graphicx}
\graphicspath{{images/}}

\newcommand{\UnitNumberName}{Unit 4: Probability, Random Variables, and Probability Distributions \quad}
\newcommand{\LessonNumberName}{Lesson 4.2: Estimating Probabilities Using Simulation}

\begin{document}
\begin{center}
    {\Large\bfseries \CourseName: \SchoolYear \\
    {\small\bfseries \UnitNumberName \LessonNumberName}}
\end{center}

\begin{tcolorbox}[colback=sky, colframe=navy, boxrule=0.9pt, arc=2mm,
    left=2mm, right=2mm, top=1.5mm, bottom=1.5mm]
    \textbf{Primary Objective:} Students will use simulation to estimate the probability
    of an event, record relative frequencies, and observe the law of large numbers in action
    (Big Idea: UNC-2).
\end{tcolorbox}

\renewcommand{\arraystretch}{1.6}
\begin{skillbox}[Priority Ideas \& Skills]{goldbox}
\begin{minipage}[t]{0.48\linewidth}
    \textbf{AP Skill 3 --- Using Probability and Simulation}
    \begin{itemize}
        \item Estimate probabilities using simulation (Skill 3.A).
        \item UNC-2.A: describe and run a simulation of a random process.
        \item UNC-2.A.4--2.A.6: design simulation (assign values, run trials, record
              outcomes) and use relative frequency to estimate probability.
        \item Apply the law of large numbers (UNC-2.A.6).
    \end{itemize}
\end{minipage}
\hfill
\begin{minipage}[t]{0.48\linewidth}
    \textbf{Key Understandings}
    \begin{itemize}
        \item A simulation is a model of a random process using a chance device (coins,
              dice, random number table, technology).
        \item Relative frequency of simulated outcomes estimates probability.
        \item As the number of trials increases, simulated (empirical) probabilities
              approach the true probability (Law of Large Numbers).
        \item Proper simulation design requires: identifying the random process, assigning
              outcomes to chance values, and recording results systematically.
    \end{itemize}
\end{minipage}
\end{skillbox}

\begin{skillbox}[{Vocabulary, Concepts, \& Theorems}]{greenbox}
\begin{minipage}[t]{0.48\linewidth}
    \begin{itemize}
        \item \textbf{Simulation} --- a model of a random process using a chance device
        \item \textbf{Trial} --- one run of a simulation
        \item \textbf{Outcome} --- the result of a single trial
        \item \textbf{Event} --- a collection of outcomes
    \end{itemize}
\end{minipage}
\hfill
\begin{minipage}[t]{0.48\linewidth}
    \begin{itemize}
        \item \textbf{Relative frequency} --- (number of times event occurred) $\div$ (total trials)
        \item \textbf{Empirical probability} --- probability estimated from simulation or real data
        \item \textbf{Law of Large Numbers} --- as the number of trials increases, empirical
              probability approaches the true probability
    \end{itemize}
\end{minipage}
\end{skillbox}

\begin{skillbox}[Activate Prior Knowledge \& Spiral Review]{sky}
\begin{minipage}[t]{0.48\linewidth}
    \textbf{Prior Knowledge Check (3--4 min)}
    \begin{itemize}
        \item Review: ``What pattern did we examine in Lesson 4.1?''
        \item Connect: In 4.1 we asked ``Could chance produce this?'' --- today we build
              a tool to actually answer that question.
        \item Quick review: relative frequency from Unit 1 (frequency tables).
        \item Prompt: \textit{``If I want to estimate how often a thumbtack lands point-up,
              what would I do?''}
    \end{itemize}
\end{minipage}
\hfill
\begin{minipage}[t]{0.48\linewidth}
    \textbf{Connections Forward}
    \begin{itemize}
        \item Formal probability rules $\to$ Lesson 4.3
        \item Conditional probability $\to$ Lessons 4.4--4.5
        \item Simulation as a check on probability calculations $\to$ throughout Unit 4
        \item Inference: simulation as a model for the null distribution $\to$ Unit 6
    \end{itemize}
\end{minipage}
\end{skillbox}

\begin{skillbox}[Hook \& Lesson]{sky}
\begin{minipage}[t]{0.48\linewidth}
    \textbf{Hook (5--7 min)}\\
    Hold up a thumbtack. Ask: ``If I drop this 50 times, what fraction of the time will
    it land point-up?'' Take student guesses. Then actually drop it (or simulate with a
    coin): record each outcome. Compute the running relative frequency after 5, 10, 20,
    50 drops. Plot it live.\\
    Transition: \textit{``We just ran a simulation. Let's formalize what we did.''}
\end{minipage}
\hfill
\begin{minipage}[t]{0.48\linewidth}
    \textbf{Lesson Flow (15--18 min)}
    \begin{enumerate}
        \item Define simulation; contrast with theoretical probability.
        \item Simulation design framework: (a) identify the random process, (b) assign
              values to outcomes, (c) use a chance device, (d) record relative frequency.
        \item Work through example: a genetic carrier cross (each parent carries a
              recessive gene). Use a coin to model each parent's allele. Two-coin trial
              = one offspring.
        \item Compute relative frequency after 10 and 30 trials; compare class results.
        \item Law of Large Numbers: show a graph of relative frequency vs.\ number of
              trials converging toward the true probability.
    \end{enumerate}
\end{minipage}
\end{skillbox}

\begin{skillbox}[Explicit Instruction \& Active Monitoring]{sky}
\begin{minipage}[t]{0.48\linewidth}
    \textbf{Direct Instruction (10 min)}
    \begin{itemize}
        \item Project: ``A free-throw shooter makes 72\% of shots. Use a random number
              table (digits 00--99) to simulate 20 shots: digits 00--71 = make, 72--99 =
              miss. Estimate the probability of making 3 in a row.''
        \item Walk through line-by-line reading of the random number table.
        \item Record outcomes in a table; compute relative frequency.
    \end{itemize}
\end{minipage}
\hfill
\begin{minipage}[t]{0.48\linewidth}
    \textbf{Active Monitoring}
    \begin{itemize}
        \item Listen for: confusing ``outcome'' with ``event'' --- remind that an event is
              a collection of outcomes.
        \item Circulate: check that students are correctly identifying which digits
              represent which outcomes.
        \item Cold-call: ``How many trials would you need before you trust your simulated
              probability? Why?''
    \end{itemize}
\end{minipage}
\end{skillbox}

\begin{skillbox}[Group Work \& Differentiation]{redbox}
\begin{minipage}[t]{0.48\linewidth}
    \textbf{Partner Activity (8--10 min)}\\
    Three-tier simulation activity: design and run a simulation for a genetics cross
    scenario (Tier R), a multi-event sports scenario (Tier A), or a custom design task
    (Tier E). Record results, compute relative frequencies, and compare to the true
    probability if known.
\end{minipage}
\hfill
\begin{minipage}[t]{0.48\linewidth}
    \textbf{Differentiation}
    \begin{itemize}
        \item \textit{Support (Tier R)}: Provide a pre-filled outcome table; students
              only need to run trials and record.
        \item \textit{On-grade (Tier A)}: Students design and run the full simulation.
        \item \textit{Extension (Tier E)}: Students run the simulation with 50 trials
              and explain what happens to the estimate as trials increase.
        \item \textit{ELL}: Vocabulary card with ``trial,'' ``outcome,'' and ``relative
              frequency'' with visual examples.
    \end{itemize}
\end{minipage}
\end{skillbox}

\begin{skillbox}[Individual Work \& Assessment]{redbox}
\begin{minipage}[t]{0.48\linewidth}
    \textbf{Exit Ticket (5 min)}\\
    A bag contains 3 red and 7 blue marbles. A student draws one marble, records the
    color, and replaces it. She repeats this 10 times and gets red 4 times.
    \begin{enumerate}
        \item What is the empirical probability of drawing red from her simulation?
        \item What is the true probability of drawing red?
        \item Does the law of large numbers predict these will get closer as she does
              more trials? Explain.
    \end{enumerate}
\end{minipage}
\hfill
\begin{minipage}[t]{0.48\linewidth}
    \textbf{AP-Style Check}\\
    \textit{``A simulation is run 200 times to estimate the probability of getting at
    least 2 heads in 3 coin flips. The event occurs in 146 of the 200 trials. Which
    of the following is the best estimate of the probability?''}
    \begin{enumerate}[label=(\Alph*)]
        \item $200/146$
        \item $146/200 = 0.73$
        \item $3/8$
        \item $5/8$
    \end{enumerate}
    Answer: (B) --- the empirical probability is $146/200 = 0.73$ (true value $= 0.5$;
    more trials would make this estimate more accurate).
\end{minipage}
\end{skillbox}

\begin{skillbox}[Reinforcement \& Extension]{goldbox}
\begin{minipage}[t]{0.48\linewidth}
    \textbf{Homework / Practice}
    \begin{itemize}
        \item Design a simulation for a real-world scenario (health screening, genetics,
              sports), run at least 20 trials, record outcomes, and compute relative
              frequency.
        \item Explain how the law of large numbers would affect the estimate if the number
              of trials were increased to 500.
    \end{itemize}
\end{minipage}
\hfill
\begin{minipage}[t]{0.48\linewidth}
    \textbf{Extension / Preview}
    \begin{itemize}
        \item \textit{Extension}: Use a random number generator (calculator or app) to
              run 100 trials and compare to the 20-trial estimate. Describe what you notice.
        \item \textit{Preview}: Lesson 4.3 introduces formal probability rules --- the
              true (theoretical) probabilities that simulation approximates. Think about what
              makes a probability ``true'' rather than estimated.
    \end{itemize}
\end{minipage}
\end{skillbox}

\end{document}
